\emph{(1).} This restates clause (G) of the premise, which says the one-slot space is a single doublet; we give the representation-theoretic form. Under clause (a) of Premise \claimref{standard-model.slot-module-hypothesis}, which is the setting of Theorem \claimref{synthesis.cpt-closure-forces-slot-module-and-chirality}, $F$ is a finite-dimensional polynomial representation of $U(2)$ graded so that $\mu I_{2}$ acts on $F_{n}$ as $\mu^{n}$. A polynomial representation is a sum of polynomial irreducibles (complete reducibility), and these are labelled by highest weights $(\lambda_{1}, \lambda_{2})$, $\lambda_{1} \ge \lambda_{2} \ge 0$, with torus weights $(p, q)$, $p + q = \lambda_{1} + \lambda_{2}$, running over $q = \lambda_{2}, \dots, \lambda_{1}$, so of dimension $\lambda_{1} - \lambda_{2} + 1$, and with $\mu I_{2}$ acting as $\mu^{\lambda_{1} + \lambda_{2}}$ (standard). On $F_{1}$ the degree is one, so every irreducible summand is $(1, 0)$, the doublet, of dimension two: $\dim F_{1}$ is even. If $F$ is generated in the sense of (G), as in that theorem, $\dim F_{0} = 1$ and an equivariant map $\C^{2} \otimes F_{0} \to F_{1}$ is onto, so $\dim F_{1} \le 2$, and $F_{1}$ is the doublet or zero (with CPT closure, Theorem \claimref{synthesis.cpt-closure-forces-slot-module-and-chirality}(3) gives $F \cong \Lambda^{\bullet}\C^{2}$, so the doublet). A three-dimensional one-slot space is excluded whatever the action of $U(2)$. The script enumerates the irreducibles of degree $0$ to $4$ with their weights and dimensions and confirms that degree one has only dimension two.

\emph{(2).} Take the face $F$ with remaining slot $s_{o}$, the gauge group acting by $U(2)$ on $S_{F}$, trivially on $s_{o}$, and by $U(3)$ on $W$. If the slot number counts all three slots, the state $s_{o}$ has $n = 1$ while $\mu I_{2}$ acts on it as $1 = \mu^{0}$, which clause (a) forbids; the same holds on the $32$ monomials containing $s_{o}$. If the slot number counts only the face's slots, then $F = \Lambda^{\bullet}S_{F} \otimes \Lambda^{\bullet}(\C s_{o})$ and $F_{0} = \mathrm{span}\{1, s_{o}\}$, against $\dim F_{0} = 1$; the map $x \mapsto s_{o} \wedge x$ is an injective equivariant map from the monomials without $s_{o}$ onto those with it, so every weight (weak $T_{3}$, face slot number, colour occupation) occurs exactly twice. This is the signature the premise names in its refutation condition: more than one slot vacuum, showing itself as two copies of each species. The parity of the total mode number differs between $x$ and $s_{o} \wedge x$, so the even half consists of the face's even states (its particle half) and the images of its odd states, which carry the weights of its antiparticle half: a generation together with its conjugate, a vectorlike pair. The script checks all of this for each of the three faces on the $64$ exterior states, including the rank $16$ of the image of the odd face states under $x \mapsto s_{o} \wedge x$ (GA's exterior model). Now let one gauge group act. By Theorem \claimref{synthesis.faces-of-one-slot-space-are-not-three-generations-under-one-gauge-group}(1), at most one face plane is an irreducible invariant subspace for any algebra acting on the slots, so at most one face's states form a slot module for that group, and by the above the shared space does not either. If each of three generations is to be a composite satisfying the premise for the one gauge group, the three state spaces are therefore three separate slot modules $\Lambda^{\bullet}(\C^{2} \oplus W_{F})$, each with its own slots and, by clause (a) (one pair system per composite), its own pair modes; the faces of Theorem \claimref{families.three-identical-families-mixing-shared-slots} enter only as labels. This is the conclusion of Theorem \claimref{synthesis.faces-of-one-slot-space-are-not-three-generations-under-one-gauge-group}(3), reached there through the overlap counts and the gauge-group lemma; the present step adds that the shared space fails the premise for one face already, through the premise's own refutation signature.

\emph{(3).} Every clause of the premise speaks of one composite's state space; a set of $n$ separate composites satisfies it composite by composite. The particle half of one slot module (chirality $\pm(-1)^{n + k}$, the even part, by the premise's derivation) has the $16$ states $(n, k) = (0,0), (0,2), (1,1), (1,3), (2,0), (2,2)$ with hypercharges $0, -\tfrac23, \tfrac16, -\tfrac12, 1, \tfrac13$ and multiplicities $1, 3, 6, 2, 1, 3$. The four conditions of Fact \claimref{standard-model.imported-anomaly-coefficients} are sums over states that are linear in the content: $\sum Y$, $\sum Y^{3}$, $\sum Y T_{3}^{2}$ (weak), $\sum Y h^{2}$ (colour, $h$ a colour Cartan generator), and the colour cubic $\sum h_{8}^{3}$; each vanishes for one copy (the script computes them exactly from the states) and so equals $n \cdot 0$ for $n$ copies. A copy has four weak doublets (three colours of the quark doublet and the lepton doublet), so $n$ copies have $4n$, an even number for every $n$, which is the condition of Witten's global $SU(2)$ anomaly (Fact \claimref{synthesis.imported-witten-global-su2-anomaly}). Nothing in the premise or in anomaly freedom refers to $n$, so they hold for every $n \ge 1$, and the number of generations is independent of them. This is standard per-generation anomaly cancellation, read as an independence statement. The script checks $n = 1, \dots, 10$ and that the wrong normalisation $Y = (3n - k)/6$ is anomalous.

\emph{(4).} Theorem \claimref{synthesis.one-loop-running-coefficients-from-spacetime-algebra-spin-traces}(4), under its Krein premise (\claimref{synthesis.krein-form-premise-indefinite-inner-product-for-gauge-fields}), gives $b_{3} = -11 + \tfrac23\sum T$ and $b_{2} = -\tfrac{22}{3} + \tfrac23\sum T + \tfrac{n_{h}}{6}$ with the Weyl Dynkin sums of the content. Computed from the states, $\sum T_{3}^{2} = 2$ and $\sum h^{2} = 2$ per copy ($h = \mathrm{diag}(\tfrac12, -\tfrac12, 0)$), giving the stored formula. $b_{3} < 0$ exactly when $\tfrac43 n < 11$, $n \le 8$; $b_{2} < 0$ exactly when $\tfrac43 n < \tfrac{22}{3} - \tfrac{n_{h}}{6}$, which for $n_{h} = 1$ and $2$ is $n \le 5$. Both coefficients increase with $n$, so the bounds hold for all $n$; the script checks $n \le 12$. This is the stored condition $8n_{g} + n_{h} < 44$ of Theorem \claimref{synthesis.one-loop-running-coefficients-from-spacetime-algebra-spin-traces}(4) at $n_{h} = 1, 2$, standard; it bounds $n$ only if both non-abelian factors are required to be asymptotically free, which no stored result requires of the weak factor.

\emph{(5).} In $\Cl(3,0)$ the grade-four part of $B^{2}$ vanishes for every bivector $B$ (there is no grade four), so $B \wedge B = 0$ and every bivector is a blade, the plane of a face; the coordinate faces $e_{2}e_{3}$, $e_{3}e_{1}$, $e_{1}e_{2}$ span the three bivectors. For $X \in GL(3)$, writing $a \wedge b = I_{3}(a \times b)$, $(Xa) \times (Xb) = \operatorname{cof}(X)(a \times b)$ with $\operatorname{cof}(X) = \det(X)X^{-\mathsf T}$ (standard); so $\Lambda^{2}S \cong S^{*} \otimes \det$. The script verifies the identity on all three faces for two rational matrices, computing the wedges with GA, and that $\operatorname{cof}(X)/\det X$ is the inverse transpose. $\binom{N}{2} = N$ means $N(N - 3) = 0$. So $\Lambda^{2}S$ is the dual triplet of the slot space twisted by the determinant. The reading that the generation labels are the face space is an argued reading under a further hypothesis, not a consequence of (2): separate composites carry no shared slot and need no labels from the maximal composite at all, so taking $\Lambda^{2}S$ as their label space, with the maximal hierarchy's slot $\su(3)$ (Theorem \claimref{standard-model.gauge-algebra-minimal-massive-composite}) acting on it, is a new family-symmetry hypothesis. Under it the count is $3$, and it rests in addition on the three-spinor bound, imported at observation level (Fact \claimref{families.imported-face-valued-bracket-quantum-numbers}). The result is transport-absent: it is representation theory on the slot module.
